\scanpage{16}
en $z$, il existe une partie ouverte \deleted{\unsure{de}} \inserted{non vide} $U$ de $Z$ telle que $z'\in U$ implique que $\unsure{\eta|y}$ est \deleted{étale en $z'$} modulaire en $z'$ [ou $=$ simplement \uline{étale} en $z'$]

\figure{\begin{tikzpicture}\draw[densely dotted](0,0)--(1.4,0);\end{tikzpicture}}

\uline{Corollaire 1} Pour que $F$ soit représentable par un $X$ \uline{étale} et \uline{séparé} sur $S$, il f et s qu'il satisfasse les conditions \deleted{1°) à 6°)} \inserted{1° 2° 3° 4° 6°}, et de plus

\deleted{\unsure{(Et)}} Le foncteur $F$ est étale (ou seulement, étale pour des arguments artiniens) \deleted{les $B$ de 5°) \unsure{sont} étales sur $A$}

\margin{\uline{Expliciter aussi le cas d'une immersion ouverte} [dont \unsure{le axiome|l'axiome} \unsure{direct|discret} est plus faible]}

\uline{Corollaire 2} Pour que $F$ soit représentable par un $X$ non ramifié et séparé sur $S$, il f. et s que les conditions 1°) à 6°), \deleted{\unsure{soient}} \inserted{ainsi que \textcircled{8°}} soient vérifiées, et de plus la condition \deleted{7} (NR.) \inserted{(qui implique 7°)}

\begin{enumerate}
\item[a)] Le foncteur $F$ est non ramifié (ou seulement, les $B$ dans 5°) sont non ramifiés sur $A$
\item[\textcircled{b}] Pour $\forall$ \deleted{\unsure{un point gén.}} \inserted{anneau local noeth} \deleted{$\unsure{A}$} \inserted{complet} de dim 1 \deleted{\unsure{d'un point de localité}} de $S$, \deleted{\unsure{avec $A$ premier}}, \inserted{irréductible}, d'idéal \deleted{et $\forall$ \unsure{$\eta\in$|$y\in$}} \inserted{premier minimal} $\mathfrak p$, et $\forall$ \deleted{$\unsure{\eta\in F|y\in F}$} \inserted{idéal} $I\subset\mathfrak p$, \inserted{tel que $I\mathfrak p=0$}, $\unsure{\eta|y}\in F(A/I)$, si $\unsure{\eta|y}_{\mathfrak p}\in\operatorname{Im}F(A_{\mathfrak p})$ alors $\unsure{\eta|y}\in\operatorname{Im}F(A)$.
\end{enumerate}

\margin{[Exprimer 8° ici par une impossibilité \uline{générique} de prolonger \ldots}

\uline{Remarque 3} La condition 8°) est conséquence de la \unsure{repré.} $\Rightarrow$ \unsure{Ce qu'il est néc.|Ce qui est néc.} pour que la conclusion soit \unsure{val. (loc.)}:
\mspage{15}
